\section{Nükleer Kuvvetlerin Vakum Doyumu ile Geometrik Birleşimi}
\label{sec:nukleer_kuvvetler}

Geleneksel fizikte kütleçekim, elektromanyetizma, güçlü nükleer kuvvet ve zayıf nükleer kuvvet dört ayrı ve uzlaştırılamaz etkileşim olarak kabul edilir. Birleşik Büzülme Hipotezinde ise tüm bu kuvvetler, \textbf{farklı mesafe ve ölçeklerdeki vakum akış geometrisinin} doğal tezahürleridir.

\subsection{Güçlü Nükleer Kuvvet ve Kuark Hapsi (Vakum Doyum Limiti)}
Makroskopik mesafelerde kütleçekimsel vakum akışı küresel olarak yayılır: $v(r) = -\sqrt{2GM/r}$ ve $g(r) \propto 1/r^2$. Bu rejimde vakum dokusu serbesttir ve akış doyuma ulaşmamıştır.

Ancak femtometre ölçeklerine ($r \sim 10^{-15}\text{ m}$) inildiğinde, kuarkların aşırı yüksek büzülme yoğunluğu nedeniyle çekilen vakum akışı \textbf{vakumun maksimum taşıma kapasitesine (Planckian doyum sınırına)} ulaşır. Vakum akısı daha fazla küresel olarak yayılamaz ($1/r^2$ seyrelemez); akış iki kuark arasında 1-boyutlu silindirik bir girdap tüpüne (\textit{flux tube} / vakum sicimi) sıkışır.

\begin{enumerate}
    \item \textbf{Kuark Hapsi (Quark Confinement):}
          Vakum tüpünün kesiti doyuma ulaştığından, tüpün birim uzunluk başına enerjisi sabittir: $\sigma \approx 1\text{ GeV/fm} \approx 1.6 \times 10^5\text{ N}$ (sicim gerilimi). İki kuark arasındaki toplam potansiyel:
          \begin{equation}
              V_{\mathrm{guclu}}(r) = -\frac{4}{3}\frac{\alpha_s \hbar c}{r} + \sigma r.
              \label{eq:guclu_potansiyel}
          \end{equation}
          Kuarkları birbirinden ayırmaya çalıştığınızda aradaki mesafe $r$ arttıkça potansiyel enerji $\sigma r$ lineer olarak sonsuza gider. Gerilim o kadar yükselir ki vakum dokusu yırtılır ve boşluktan yeni bir kuark-antikuark çifti ($q\bar{q}$) doğar; bu nedenle evrende asla serbest tek bir kuark gözlemlenemez!

    \item \textbf{Asimptotik Serbestlik (Asymptotic Freedom):}
          Kuarklar birbirine çok yaklaştığında ($r \to 0$), büzülme havzaları üst üste biner. Aradaki vakum doyum tüpü gevşer ve lineer terim sıfırlanır ($\sigma r \to 0$). Kuarklar çekirdek içinde neredeyse tamamen serbest parçacıklar gibi davranır.

    \item \textbf{Topolojik Renk Yükü ve $120^\circ$ Kiral Faz Dengesi ($SU(3)$ Çözümü):}
          Standart Modeldeki üç soyut renk yükü ($SU(3)$ singleti / ``renksiz hadron'' şartı), Ölçek-Zaman Kalkülüsünde üç kuark girdabının ($\mathbb{S}_k = \mathbf{X}_k + \kappa \lambda e^{i \phi_k}, \kappa^2 = 0$) dış uzaya net bir dönel faz burulması sızdırmadan en kararlı kapalı havzayı oluşturabilmesi için gereken \textbf{$120^\circ$ ($2\pi/3$) simetrik kiral faz dengesidir}:
          \begin{equation}
              \sum_{k=1}^3 e^{i \phi_k} = 1 + e^{i 2\pi/3} + e^{i 4\pi/3} = 1 + \left( -\frac{1}{2} + i\frac{\sqrt{3}}{2} \right) + \left( -\frac{1}{2} - i\frac{\sqrt{3}}{2} \right) \equiv 0.
              \label{eq:renk_faz_toplami_main}
          \end{equation}
          Tıpkı üç fazlı dengeli bir alternatif akım devresinde nötr hattaki toplam akımın sıfırlanması gibi, $120^\circ$ faz farkıyla dönen üç girdap hadronun dış sınırında birbirinin teğetsel girdap sürüklenmesini tam olarak iptal eder. Dış uzaya net bir burgu sızmaz; hadron dışarıdan bakıldığında bütünüyle nötr ve ``renksiz'' görünür. Mezonlarda ($q\bar{q}$) ise bu denge zıt kiralite ve nilpotent süperpozisyonla ($\kappa^2 = 0$) sağlanır.
\end{enumerate}

\subsection{Zayıf Nükleer Kuvvet: Girdap Sızıntısı ve Kiral Faz Kırınımı}
Zayıf nükleer etkileşim (örneğin nötronun beta bozunumuyla proton, elektron ve antinötrinoya dönüşmesi: $n \to p + e^- + \bar{\nu}_e$), büzüşen nükleon girdabındaki \textbf{kiral faz sızıntısıdır}.

Büzülme süreci sırasında kuark girdabının içsel fazı, statik vakumun aşırı yüksek doyum basıncı altında küçük bir faz kaymasına (\textit{chirality shear}) uğrar. Bu sızıntı, vakum üzerinde yerel bir faz dalgalanması olarak $W^\pm$ veya $Z^0$ kütleli vektör bozonları şeklinde paketlenir. $W$ bozonunun $80\text{ GeV}$'lik devasa kütlesi, doymuş nükleer vakum duvarını yırtabilmek için gereken minimum faz gerilim eşiğidir. Zayıf etkileşimin son derece kısa menzilli ($r_W \approx \hbar / (M_W c) \sim 10^{-18}\text{ m}$) olması, bu sert doyum bariyerinden kaynaklanan geometrik bir bastırmadır (Fermi sabiti $G_F \sim 1/M_W^2$).

\begin{enumerate}
    \item \textbf{Serbest Nötron Bozunumu ($N=1$):} Yalıtılmış serbest bir nötronda dış vakum kalkanı bulunmadığından, Ölçek-Kovaryant Türev $\mathbb{D}_t = \frac{\partial}{\partial t} + \kappa H_0$ altındaki içsel kiral faz burulması WKB tünellemesiyle doyum bariyerini aşar ($\tau_n \approx 879\text{ s} \approx 14.6\text{ dakika}$).
    \item \textbf{Çekirdek İçi Kararlılık (Kolektif Faz Kilitlemesi):} Nötron bir atom çekirdeğine girdiğinde komşu nükleonların akı tüpleri ortak bir kalkan örer ($\mathbb{S}_{\mathrm{cekirdek}} = \sum \mathbb{S}_i$). Çekirdeğin nükleer bağlanma enerjisi ($\Delta B \sim 8\text{ MeV/nükleon}$), beta bozunumunun $Q \approx 0.782\text{ MeV}$ marjını 10 kat aştığından nilpotent projeksiyon operatörü $\hat{\mathcal{P}}_\kappa$ sızıntı kanalını spektral olarak kilitler. Serbestken 15 dakikada bozunan nötron, çekirdek içinde evrenin ömründen daha uzun süre ($> 10^{30}\text{ yıl}$) mutlak bir kararlılık kazanır.
\end{enumerate}

\begin{quote}
    \small
    \textbf{Kavramsal ve Fiziksel Özet (Kuvvetlerin Efektif Geometrik Rejimleri):} \\
    1. \textbf{Kütleçekim ($\gamma = 1$ Rejimi):} Asimptotik makroskopik mesafelerde küresel seyrelti altındaki serbest vakum akış potansiyelidir ($v^2 \propto 1/r$, $\Phi \propto 1/r$). \\
    2. \textbf{Güçlü Etkileşim ($\gamma = -1$ Rejimi):} Femtometre ölçeğinde vakum taşıma kapasitesinin doyum sınırına ulaşmasıyla oluşan 1-boyutlu silindirik akı tüpü gerilimidir ($V(r) \propto \sigma r$). Kuark hapsi ve asimptotik serbestlik, bu doyum geometrisinin topolojik sonucudur. \\
    3. \textbf{Zayıf Etkileşim (Kiral Tünelleme Rejimi):} Doymuş nükleer vakum sınırından kaynaklanan kiral faz kayması ve WKB faz sızıntısıdır; $W^\pm$ ve $Z^0$ bozonlarının etkin kütleleri bu doyum bariyerinin aşılması için gereken minimum faz gerilimidir.
\end{quote}

\subsection{Nötrino Salınımları ve PMNS Matrisinin Kiral Faz Rezonansı}
\label{sec:notrino_pmns_main}

Standart parçacık fiziğinde kuark sektöründeki Cabibbo--Kobayashi--Maskawa (CKM) karışım açıları son derece küçükken ($\theta_{12}^{\mathrm{CKM}} = \theta_C \approx 13^\circ$, $\theta_{23}^{\mathrm{CKM}} \approx 2.4^\circ$, $\theta_{13}^{\mathrm{CKM}} \approx 0.2^\circ$); lepton sektöründeki Pontecorvo--Maki--Nakagawa--Sakata (PMNS) nötrino karışım açılarının devasa ve maksimal olması ($\theta_{12} \approx 33.4^\circ$, $\theta_{23} \approx 45^\circ$, $\theta_{13} \approx 8.6^\circ$) temel bir simetri bilmecesidir.

Birleşik Büzülme Hipotezinde bu hiyerarşi, \textbf{hapis geometrisi ile serbest faz yayılımı arasındaki topolojik zıtlıktan} türer:
\begin{enumerate}
    \item \textbf{Kuark Sektörü (CKM / Küçük Karışım):} Kuarklar femtometre ölçeğinde 1-boyutlu doymuş silindirik akı tüpüne ($\gamma = -1$) hapsolmuştur. Bu katı doyum geometrisi, kuark girdaplarının faz kaçıklığını dar bir açısal aralığa sıkıştırır; Cabibbo açısı nükleon büzülme sınırında küçük bir teğetsel kayma olarak belirir ($\theta_C \approx 13^\circ$).
    \item \textbf{Nötrino Sektörü (PMNS / Büyük ve Maksimal Karışım):} Nötrinolar doymuş bir akı tüpüne hapsolmayan, durağan vakum dokusunda serbestçe yayılan kütlesiz veya ultra-hafif saf kiral faz girdaplarıdır. Hadronik bir sınır kısıtlaması bulunmadığından, üç lezzet modu ($\nu_e, \nu_\mu, \nu_\tau$) statik vakumun global kiral gerilimiyle makroskopik mesafeler boyunca serbest faz rezonansına girer. Lezzet durumları ile vakum kütle öz-durumları arasındaki dönüşüm:
          \begin{equation}
              |\nu_\alpha\rangle = \sum_{i=1}^3 U_{\alpha i}^{\mathrm{PMNS}} |\nu_i\rangle
              \label{eq:pmns_donusum}
          \end{equation}
          şeklindedir. İki nötrino modu arasındaki salınım olasılığı:
          \begin{equation}
              P(\nu_\alpha \to \nu_\beta) = \sin^2(2\theta) \sin^2\left( \frac{\Delta m^2 c^3 L}{4 \hbar E} \right)
              \label{eq:notrino_salinim}
          \end{equation}
          olup, vakumun kiral gerilimi altında atmosferik nötrinolar için rezonans maksimal sınıra ($\theta_{23} \approx 45^\circ$), güneş nötrinoları için ise geniş açıya ($\theta_{12} \approx 33.4^\circ$) oturur.
\end{enumerate}
Böylece kuarklardaki küçük karışım ile nötrinolardaki maksimal karışım, 1D hadronik hapis ile serbest vakum yayılımı arasındaki doğrudan geometrik fark olarak çözülür.

\subsection{Kuvvet Hiyerarşisi ve Yoğunluk Kaynaklı Boyutsal Geçiş: Atom Sayısı, Planck Doyumu ve Geometrik Boyut (\texorpdfstring{$\gamma$}{gamma})}
\label{sec:boyutsal_gecis_main}

Standart fizikte kütleçekim ile mikroskobik kuvvetler arasındaki $\sim 10^{36}$--$10^{40}$ katlık devasa şiddet uçurumu ``Hiyerarşi Problemi'' (Hierarchy Problem) olarak adlandırılır. Birleşik Büzülme modelinde bu oran, **ortamdaki madde-enerji yoğunluğunun vakumun Planck doyum eşiğine ($c^4/G$) oranına bağlı olarak gelişen boyutsal bir faz geçişidir (Density-Driven Dimensional Crossover)**:
\begin{equation}
    \left( \frac{v(r)}{c} \right)^2 = 2 \, \mathcal{C}_0 \left( \frac{r}{\ell_0} \right)^{-\gamma}.
    \label{eq:master_formula_main}
\end{equation}
Burada $\mathcal{C}_0$ boyutsuz temel kuplaj katsayısı, $\ell_0$ karakteristik geçiş ölçeğidir (örneğin kütleçekimde Schwarzschild yarıçapı $r_s$, hadronik sektörde femtometre $\ell_{\mathrm{QCD}}$); bu sayede $\gamma$ boyutsal indisi sürekli değişse dahi formül her rejimde boyutsal olarak mutlak tutarlılığını muhafaza eder.

Burada kritik soru şudur: \textit{``Kuvvetlerin şiddetini belirleyen şey atom sayısı mıdır, yoksa geometrik boyut ($\gamma$) mudur?''}

Bu iki kavram bir neden-sonuç zinciridir:
\begin{enumerate}
    \item \textbf{Güneş Örneği (Büyük $N$, Düşük Yoğunluk $\implies \gamma = 1$):} Güneş'te $N \sim 10^{57}$ atom bulunur. Devasa atom sayısı büyük bir kütleçekim debisi ($M = \sum m_i$) üretse de, bu atomlar yüz binlerce kilometrelik astronomik bir hacme yayılmıştır. Birim hacme düşen vakum akısı, vakumun Planckian doyum eşiğinin trilyonlarca kat altında kalır. Akış $4\pi r^2$ küresel kabuğu üzerinden serbestçe seyrelebildiği için geometri 3-boyutludur ($\gamma = 1$). Kütleçekimin aşırı zayıf görünmesi, debinin 3D uzayda küresel olarak seyrelmesinden kaynaklanır.
    \item \textbf{Proton Örneği (Küçük $N=3$, Ultra-Yüksek Yoğunluk $\implies \gamma = -1$):} Bir protonun içinde yalnızca 3 kuark bulunur. Atom sayısı bakımından Güneş'ten $10^{57}$ kat daha küçüktür; ancak bu 3 kuark $r \sim 10^{-15}\text{ m}$ (femtometre) gibi akıl almaz küçüklükte bir çekirdeğe hapsolmuştur. Bu ölçekte büzülme akı yoğunluğu vakum dokusunun maksimum taşıma kapasitesini ($\Phi_{\max} = c^4/G$) aşar! Vakum daha fazla küresel olarak seyrelemez ($1/r^2$ tükenir); dinamik sistem en düşük direnç yolu olarak **1-boyutlu silindirik bir girdap tüpüne ($\gamma = -1$)** çöker. Tüpün kesiti sabitlendiğinden mesafe arttıkça gerilim düşmez ($\sigma \approx 16\text{ ton}$) ve kuark hapsi doğar.
    \item \textbf{Zayıf Kuvvet (Doyum Bariyeri Tünellemesi $\implies \gamma \to \text{üstel}$):} Bu doymuş nükleer girdaptan dışarı tünelleyen kiral faz sızıntısı ($W^\pm, Z^0$), doyum duvarının kalınlığı sebebiyle üstel olarak sönümlenir ($\mathcal{C}_W \propto e^{-r/r_W}$).
\end{enumerate}
**Yerel yoğunluk Planck doyumunu aştığında akış 3D serbest küresellikten ($\gamma=1$) 1D sicim tüpüne ($\gamma=-1$) boyut değiştirmektedir.**

\subsubsection{Boyutsal İndeks $\gamma$'nın Sürekli Doyum Faz Geçişi Denklemi}
\label{sec:surekli_gamma}
Modelimizde $\gamma$ parametresi, akışkan dinamiğindeki doyum faz geçişini yöneten sürekli bir enerji yoğunluğu fonksiyonudur:
\begin{equation}
    \gamma(\rho) = \tanh\left( \ln\left[ \frac{\rho_{\mathrm{sat}}}{\rho(r)} \right] \right)
    \label{eq:gamma_surekli_gecis}
\end{equation}
burada $\rho_{\mathrm{sat}} \approx 10^{18}\text{ kg/m}^3$ hadronik nükleer doyum yoğunluğudur.
\begin{itemize}
    \item \textbf{Düşük Yoğunluk Rejimi ($\rho \ll \rho_{\mathrm{sat}}$):} Makroskopik kütleçekim ve atomik elektron bulutunda $\frac{\rho_{\mathrm{sat}}}{\rho} \gg 1$ olduğundan $\ln(\rho_{\mathrm{sat}}/\rho) \to +\infty$ ve $\tanh \to +1$'dir. Böylece $\gamma \equiv +1$ (3D serbest küresel akı) kesin olarak elde edilir.
    \item \textbf{Aşırı Yoğunluk Rejimi ($\rho \gg \rho_{\mathrm{sat}}$):} Kuark ölçeğinde ($r \le 10^{-15}\text{ m}$) yerel yoğunluk Planck doyumunu aştığından $\frac{\rho_{\mathrm{sat}}}{\rho} \ll 1$, $\ln(\rho_{\mathrm{sat}}/\rho) \to -\infty$ ve $\tanh \to -1$'dir. Böylece $\gamma \equiv -1$ (1D sıkışmış akı tüpü / kuark hapsi) pürüzsüz biçimde ortaya çıkar.
\end{itemize}

\subsubsection{Hiyerarşi Açmazının Çözümü: $G$'nin Cebirsel Sadeleşmesi ve Cornell Gerilimi}
\label{sec:hiyerarsi_monograph}
Vakum doyum gerilimi $F_{\mathrm{Planck}} = c^4/G \approx 1.21 \times 10^{44}\text{ N}$ seviyesindeyken, hadronik akı tüpünün geriliminin neden $\sigma \approx 1.6 \times 10^5\text{ N}$ ($\approx 1\text{ GeV/fm}$) olduğu sorusu; parçacığın kütlesi $m_q$ ile Planck ölçeği arasındaki boyutsuz kuantum kütleçekim kuplajı ($\alpha_G \equiv G m_q^2 / \hbar c$) üzerinden analitik olarak çözülür:
\begin{equation}
    \sigma = \alpha_G \cdot F_{\mathrm{Planck}} = \left( \frac{G m_q^2}{\hbar c} \right) \left( \frac{c^4}{G} \right) = \frac{m_q^2 c^3}{\hbar}.
    \label{eq:sigma_monograph_ispat}
\end{equation}
Burada $G$ sabiti pay ve paydadan tamamen yok olur! $m_q \approx 440\text{ MeV}$ kurucu kuark kütlesi için $\sigma \approx 1.57 \times 10^5\text{ N} \approx 0.98\text{ GeV/fm}$ değeri elde edilir; bu da 39 mertebelik hiyerarşiyi ad-hoc varsayımsız kapatır.

\subsection{Epistemolojik Zemin: Dört Temel Postulatın Tek Parametrede Eritilmesi}
\label{sec:postulatlarin_eritilmesi_sentez}

Standart fizik paradigmalarında (Genel Görelilik ve Standart Model) birleşik alan teorilerinin bir asırdır başarısızlığa uğramasının temel epistemolojik nedeni, dört kuvveti birleştirmeye çalışırken, bu kuvvetlerin tanımlandığı \textbf{zaman ($t$), konum ($\mathbf{x}$), kütle ($m$) ve ışık hızı ($c$)} parametrelerini birbirinden bağımsız ve bağlantısız birer aksiyom olarak muhafaza etmeleridir.

Doğadaki dört temel etkileşimi tek bir dinamik denklemde eritebilmek için, bu dört temel büyüklüğün kendilerinin de bağımsız postulatlar olmaktan çıkarılıp, tek bir ortak parametrenin—maddenin vakumdaki ölçek daralma dinamiğinin ($\lambda(t)$ ve $\kappa$)—farklı geometrik ve kinematik izdüşümleri olarak anlaşılması şarttır. Çabamızın ulaştığı bu ontolojik indirgeme Tablo~\ref{tab:ontolojik_indirgeme_sentez}'de özetlenmiştir.

\begin{table}[h!]
    \centering
    \small
    \begin{tabular}{@{}lll@{}}
        \toprule
        \textbf{Klasik/Standart Postulat} & \textbf{Ölçek-Zaman Kalkülüsündeki Karşılığı}                                                                                     & \textbf{Fiziksel Mekanizma}                              \\ \midrule
        \textbf{Işık Hızı ($c$)}          & $v_c = 1/\sqrt{\varepsilon_0\mu_0} = c$                                                                                           & Vakumun elastik faz iletim ve içsel kuantum tavan hızı   \\
        \textbf{Zaman ($t$)}              & $dt \equiv -\frac{1}{H_0} \frac{d\lambda}{\lambda} = -\frac{1}{H_0} d(\ln\lambda)$                                                & Atomların büzülürken tuttuğu kuantum faz günlüğü         \\
        \textbf{Konum / Uzaklık ($R$)}    & $\mathbf{R}(t) \equiv \frac{\mathbf{X}_0}{\lambda(t)} = \mathbf{X}_0 e^{H_0 t}$                                                   & Statik uzayın büzülen yerel cetvele oranı                \\
        \textbf{Kütle ($M$)}              & $M \equiv \frac{1}{4\pi G} \oint_{\partial \Omega} |(\mathbf{v}_{\mathrm{inflow}}\cdot\nabla)\mathbf{v}_{\mathrm{inflow}}| \, dA$ & Vakum akışının ivme alanının (hız gradyanının) net akısı \\ \bottomrule
    \end{tabular}
    \caption{Dört temel klasik postulatın tek bir ölçek daralma mekanizmasında ($\lambda(t)$) eritilmesi.}
    \label{tab:ontolojik_indirgeme_sentez}
\end{table}

Bu dört temel kavram tek bir mekanizmada çözüldüğünde; Kütleçekim, Elektromanyetizma, Güçlü Kuvvet ve Zayıf Kuvvet'in akış dinamikleri, aynı vakum sürekliliğindeki ($\mathbf{v}_{\mathrm{inflow}}$) farklı yoğunluk ve doyum rejimlerinin geometrik formları olarak sınıflandırılabilir.

\subsection{Geometrik Sentez: Dört Etkileşimin Birleşik Akı Dinamiği}
\label{sec:ana_dinamik_denklem_main}

Ölçek-Zaman Kalkülüsünün gözlemci koordinatındaki kinematik ivme bağıntısı ile Bütünleşik Vakum Akı Ana Formülü (\ref{eq:master_formula_main}) birleştirildiğinde, doğadaki tüm mikroskobik ve makroskobik dinamikleri tek bir çatıda toplayan \textbf{Bütünleşik Ana Dinamik Denklem} elde edilir:
\begin{equation}
    \boxed{ m \frac{d^2\mathbf{R}}{dt^2} = - m c^2 \nabla \left[ \mathcal{C}_0 \left(\frac{R}{\ell_0}\right)^{-\gamma} \right] + 2 m H_0 \mathbf{v}_{\mathrm{yerel}} + m H_0^2 \mathbf{R} }
    \label{eq:master_dynamic_equation}
\end{equation}

Bu tek diferansiyel denklem evrenin tüm rejimlerini eksiksiz açıklar:
\begin{enumerate}
    \item \textbf{Birinci Terim (Dört Temel Kuvvetin Vakum Potansiyeli):} $- m c^2 \nabla [\mathcal{C}_0 (R/\ell_0)^{-\gamma}]$ terimi, ortamın kütle-enerji yoğunluğuna göre şekillenir:
          \begin{itemize}
              \item $\gamma = 1 \implies$ Kütleçekim ($\mathcal{C}_0 = GM/(c^2 \ell_0)$) ve Elektromanyetizma ($\mathcal{C}_0 = \alpha \hbar / (m_e c \ell_0)$).
              \item $\gamma = -1 \implies$ Güçlü Nükleer Kuvvet / Kuark Hapsi ($\mathcal{C}_0 = \sigma \ell_0 / (2m_q c^2)$).
              \item $\gamma \to \text{üstel} \implies$ Zayıf Nükleer Kuvvet / Beta Bozunumu ($\mathcal{C}_0 \propto e^{-R/r_W}$).
          \end{itemize}
    \item \textbf{İkinci Terim (Ölçek Faz Sönümlemesi):} $2 m H_0 \mathbf{v}_{\mathrm{yerel}}$, büzülen referans sistemindeki hareketin hissettiği kozmik Coriolis direncidir.
    \item \textbf{Üçüncü Terim (Kozmik Kaçıklık / Karanlık Enerji Yanılsaması):} $m H_0^2 \mathbf{R}$, standart kozmolojinin harici Karanlık Enerji ($\Lambda \approx 3 H_0^2$) ile açıkladığı kozmik ivmeli genleşme yanılsamasının, büzülen koordinat kalkülüsündeki doğal matematiksel karşılığıdır.
\end{enumerate}

Böylece teori iki kademeli olarak kesin analitik çözüme kavuşturulmuştur:
\begin{itemize}
    \item \textbf{Akı ve Hız Düzeyinde Ana Formül:} $\left(\frac{v(r)}{c}\right)^2 = 2 \, \mathcal{C}_0 \left(\frac{r}{\ell_0}\right)^{-\gamma}$ (Statik Vakum Akışı).
    \item \textbf{Kuvvet ve Yörünge Düzeyinde Ana Denklem:} $m \frac{d^2\mathbf{R}}{dt^2} = - m c^2 \nabla \left[ \mathcal{C}_0 \left(\frac{R}{\ell_0}\right)^{-\gamma} \right] + 2 m H_0 \mathbf{v} + m H_0^2 \mathbf{R}$ (Ölçek-Kovaryant Dinamik Hareket).
\end{itemize}

